% Modernised reading — fonds Alexandre Grothendieck, shelfmark 35, pages 1-5.
%
% This is an INTERPRETATION, not a transcription. Everything the critical
% apparatus carried moves into footnotes. In any dispute of substance the
% transcription governs: whoever cites Grothendieck must cite batch-01.fr.tex.
%
% Pass: Opus 5 (claude-opus-5), 2026-08-16 — first pass, unchecked against the
% pages by a human. The folder was read whole; it holds one batch, pages 1-5.
%
% The folder is a plan for a seminar that was afterwards published, which sets a
% trap: it is tempting to read every line against the printed SGA 7 and report
% what "became" of it. Comparisons with the published volumes are made only where
% the plan itself carries the change — a renumbering, a name struck and replaced —
% and are marked as ours wherever they go beyond that.
\documentclass[11pt,a4paper]{article}
% Le préambule commun à toutes les transcriptions du fonds Grothendieck.
%
% Il tient en une idée : **l'incertitude se compose, elle ne se résout pas.**
% Une transcription qui lisse un mot douteux a détruit la seule chose pour
% laquelle on la faisait — un lecteur doit pouvoir savoir, à chaque endroit,
% ce qui était sur le feuillet et ce qui a été ajouté. Les six macros qui
% suivent sont donc l'appareil critique, et non de la décoration.
%
% Les mêmes commandes sont reconnues par `scripts/render.mjs`, qui produit la
% vue de lecture du site. Une macro ajoutée ici doit l'être aussi là-bas,
% faute de quoi le rendu échouera bruyamment — ce qui est voulu : un rendu
% silencieusement faux est pire qu'un rendu absent.

\ProvidesPackage{grothendieck}[2026/08/08 Transcriptions du fonds Grothendieck]

\usepackage{amsmath,amssymb,amsthm}
\usepackage{mathrsfs}
\usepackage{stmaryrd}
% Les diagrammes commutatifs sont partout dans ce fonds : c'est la langue
% même de la géométrie algébrique de Grothendieck, et une transcription qui
% les rendrait en prose ne transcrirait rien.
\usepackage{tikz-cd}
% Français par défaut — c'est la langue des feuillets — mais l'anglais est
% chargé aussi : la traduction et le résumé sont des documents anglais, et la
% typographie française (espaces avant les deux-points) y serait une faute.
% Ces documents basculent par \selectlanguage{english} en tête de corps.
\usepackage[english,french]{babel}
\usepackage{fontspec}
\usepackage{geometry}
\geometry{a4paper,margin=25mm,marginparwidth=28mm}
\usepackage{marginnote}
\usepackage{xcolor}
% ulem, not soul. `soul`'s \st and \ul are built on a character-by-character
% scan that XeLaTeX's font machinery breaks: the document fails with an
% undefined control sequence rather than degrading. ulem does the same two jobs
% and is Unicode-engine safe. `normalem` stops it hijacking \emph, which the
% transcriptions use for Grothendieck's own emphasis.
\usepackage[normalem]{ulem}
\usepackage{hyperref}
\hypersetup{colorlinks=true,linkcolor=black,urlcolor=black,pdfborder={0 0 0}}

% --- Métadonnées du lot -----------------------------------------------------
% Reprises telles quelles par le script de rendu, qui les affiche en tête de
% la vue de lecture. Elles fixent ce que le fichier prétend couvrir : sans
% elles, une transcription flotte sans point d'attache dans un fonds de seize
% mille feuillets.

\newcommand{\folder}[1]{\gdef\grothFolder{#1}}
\newcommand{\batch}[1]{\gdef\grothBatch{#1}}
\newcommand{\leaves}[2]{\gdef\grothFirst{#1}\gdef\grothLast{#2}}
\newcommand{\foldertitle}[1]{\gdef\grothFoldertitle{#1}}
\gdef\grothFolder{?}\gdef\grothBatch{?}\gdef\grothFirst{?}\gdef\grothLast{?}\gdef\grothFoldertitle{}

% --- Le résumé --------------------------------------------------------------
%
% Il ouvre la lecture modernisée, sous le titre « Résumé », et il est composé
% en petit. Ce n'est pas un ornement : c'est le seul passage du document qui
% ne soit pas des mathématiques, et un lecteur doit pouvoir le reconnaître
% avant de l'avoir lu — puis le sauter s'il connaît déjà le sujet, ou s'y
% arrêter s'il ne le connaît pas. Le filet qui le suit marque la reprise du
% corps.

\newenvironment{resume}
  {\small\setlength{\parskip}{0.45em}}
  {\par\medskip\hrule\medskip}

% --- Mots-clés ---------------------------------------------------------------
%
% En anglais, en clôture du résumé de la lecture modernisée : le vocabulaire
% moderne sous lequel chercher ce que les feuillets construisent. Le manifeste
% du site (`npm run manifest`) les extrait pour étiqueter la cote entière —
% cette ligne est la source unique des tags, il n'y a pas de fichier à côté.
\newcommand{\keywords}[1]{\par\smallskip\noindent
  {\footnotesize\textsc{Keywords} --- \textit{#1}}\par}

% --- L'appareil critique ----------------------------------------------------

% Le numéro de feuillet, en marge. C'est l'ancre qui permet de régler un
% désaccord sur une formule en regardant *un* feuillet plutôt qu'une chemise
% de six cents. Le site s'en sert aussi pour tourner les pages du fac-similé
% quand on fait défiler la transcription.
\newcommand{\leaf}[1]{%
  \marginnote{\footnotesize\color{gray!70}\textsf{#1}}%
  \label{leaf:#1}\ignorespaces}

% Illisible : on ne devine pas. Un mot inventé qui se lit comme les autres est
% le pire résultat possible.
\newcommand{\ill}{\textcolor{red!60!black}{[\,\ldots\,]}}

% Lecture douteuse : le mot est donné, et signalé comme incertain.
\newcommand{\uncertain}[1]{\uline{#1}}

% Ajout de l'éditeur — une lettre manquante, un mot sous-entendu.
\newcommand{\add}[1]{\textcolor{blue!50!black}{[#1]}}

% Biffé par Grothendieck lui-même. Ce qu'il a rayé fait partie du document :
% une rature dit souvent où la pensée a changé de direction.
\newcommand{\struck}[1]{\sout{#1}}

% Note du transcripteur, sur l'état matériel du feuillet.
\newcommand{\note}[1]{\par\smallskip{\small\color{gray!60!black}\textit{[#1]}}\par\smallskip}

% Note marginale de Grothendieck, distincte du corps du texte.
\newcommand{\marginal}[1]{\par\smallskip\noindent\hspace{1em}%
  {\small\color{gray!60!black}#1}\par\smallskip}

% --- Titre ------------------------------------------------------------------

\AtBeginDocument{%
  \begin{center}
    {\large\bfseries Fonds Alexandre Grothendieck --- cote n\textsuperscript{o}~\grothFolder}\\[2pt]
    {\itshape\grothFoldertitle}\\[4pt]
    {\small feuillets \grothFirst--\grothLast\ (lot \grothBatch)}\\[2pt]
    {\footnotesize\color{gray!60!black}Transcription établie d'après le fac-similé numérisé
     par l'Université de Montpellier.\\
     Les lectures incertaines sont soulignées, les illisibles notées [\,\ldots\,],
     les ajouts entre crochets.}
  \end{center}
  \vspace{1em}}

\endinput


\folder{35}
\pages{1}{5}
\dating{[vers 1967-1973]}
\watermark{Édition de démonstration\\interprétation personnelle de l'œuvre}
\foldertitle{Plan [SGA 7] : notes manuscrites (s.d.), tapuscrit annoté (s.d.) — lecture modernisée du dossier entier}

\begin{document}

\section*{Résumé}

\begin{resume}
Cinq feuillets qui montrent comment on organise un séminaire de recherche : la
table des exposés, deux fois ; la bibliographie qu'on avait commencée ; un
feuillet de rédaction ; et, pour finir, deux graphes où l'on a dessiné qui
dépend de qui.

Le séminaire en question porte sur la \emph{monodromie}. Le mot désigne un
phénomène très simple à décrire et très difficile à contrôler. Prenons une
famille de courbes qui varie continûment --- disons une famille de courbes
paramétrée par un disque --- et supposons qu'une seule d'entre elles, celle du
centre, soit singulière, pincée en un point. Si l'on fait tourner le paramètre
une fois autour du centre, chaque courbe non singulière revient sur elle-même ;
mais rien n'oblige les cycles qu'elle porte à revenir à leur place. Ils
peuvent avoir été permutés, ou mêlés. Cette transformation, qu'un tour complet
inflige à l'homologie, est la monodromie ; et l'on veut savoir de quoi elle est
capable.

Deux réponses sont visées, et les deux graphes du dernier feuillet en portent
les noms. Une réponse \emph{qualitative} : la monodromie est toujours
quasi-unipotente, c'est-à-dire qu'une de ses puissances agit de façon aussi peu
compliquée que possible. Une réponse \emph{quantitative} : dans les bons cas on
sait dire précisément ce qu'elle fait, et la formule de Picard--Lefschetz le
dit.

Ce que ce dossier apporte n'est donc pas un théorème, c'est une architecture.
On y voit un problème découpé en une vingtaine de morceaux, chaque morceau
attribué à quelqu'un --- les noms sont écrits dans la marge, et plusieurs sont
raturés puis remplacés ---, et l'on y voit surtout deux tentatives successives
de dessiner l'ordre logique de ces morceaux : la première barrée de quatre
grands traits, la seconde tracée en dessous. C'est une image assez rare de ce
qu'est réellement la construction d'une théorie : pas une suite d'énoncés, mais
un ordre à trouver.

\keywords{monodromy, vanishing cycles, Picard-Lefschetz formula, semi-stable reduction, biextension, Lefschetz pencil}
\end{resume}

\pagerange{1}{1}
\section*{La première table : les exposés I à IX}

Le premier feuillet est un tapuscrit --- la table avait donc déjà été mise au
propre --- repris ensuite à l'encre. Les corrections sont de deux sortes : des
noms ajoutés en marge, un par ligne, et des retouches dans le corps.

L'architecture qui s'y lit est celle-ci.

\begin{itemize}
\item \textbf{I à III --- le théorème de monodromie, deux fois.} L'exposé I
établit les généralités sur le groupe de monodromie locale et démontre le
théorème par voie \emph{arithmétique} ; l'exposé III le redémontre par voie
\emph{géométrique}. L'exposé II traite le cas d'actions résolubles du groupe
fondamental et se présente comme une variante du même théorème.\footnote{Deux
démonstrations d'un même énoncé, données comme deux exposés distincts, ne sont
pas une redondance de plan : la première passe par l'arithmétique des corps
locaux, la seconde par la géométrie de la réduction, et ce que le séminaire
veut est précisément que les deux soient disponibles. Cette lecture est la
nôtre ; la table ne la commente pas.}

\item \textbf{IV --- les cycles évanescents.} Propriétés générales des groupes
de cycles évanescents, et le cas des points singuliers isolés de la fibre
spéciale. C'est l'objet technique central de tout le séminaire : ce qui reste
d'un cycle quand on le fait dégénérer vers la fibre singulière.

\item \textbf{V --- l'action sur le groupe fondamental.} La monodromie n'agit
pas seulement sur l'homologie mais sur le $\pi_{1}$ lui-même, et l'on veut là
des théorèmes d'action modérée et essentiellement modérée.

\item \textbf{VI à VIII --- les biextensions.} Après un exposé sur la
présentation finie des groupes fondamentaux géométriques et un autre sur les
modules formels, viennent les biextensions de faisceaux de groupes, puis les
extensions et biextensions de schémas en groupes, avec un théorème
d'orthogonalité.\footnote{Le titre de l'exposé VIII est en partie ajouté à la
main par-dessus le tapuscrit et une partie en reste illisible ; seul subsiste
clairement le $\mathbb{G}_{m}$ et le mot « orthogonalité ». La biextension est
la structure algébrique qui porte l'accouplement de dualité entre une variété
abélienne et sa duale ; c'est cette théorie qui rend possible l'exposé IX.}

\item \textbf{IX --- l'application.} Les modèles de Néron, et le théorème de
réduction semi-stable.
\end{itemize}

Un dixième exposé --- « la réduction semi-stable des courbes algébriques, et la
connexité des variétés de modules pour les courbes de genre $g$ » --- est
entièrement barré, et le nom porté en face de lui est celui de Deligne.

\subsection*{Ce que les marges disent}

La colonne des noms ajoutée à gauche est le document dans le document. La
plupart des lignes portent \emph{Gr.} --- l'auteur lui-même ---, ce qui indique
que le séminaire était, dans cet état du plan, largement le sien.

Trois lignes échappent à cette règle et méritent d'être relevées. En face de
l'exposé V, cerclé et suivi d'un point d'interrogation : \emph{Mme Raynaud ?}
--- une attribution encore hypothétique. En face de l'exposé renuméroté, un nom
écrit par-dessus un autre qui est barré, tous deux Raynaud.\footnote{Le prénom
ou la civilité du nom supérieur est peu lisible ; la transcription le marque
incertain. Ce qui est certain est qu'un Raynaud remplace un autre Raynaud, ce
qui est en soi un renseignement sur la répartition du travail.} Et en face de
l'exposé VI sur les modules formels, un \emph{Rim} barré --- l'exposé change de
main --- suivi, dans le titre lui-même, de trois mots entre parenthèses :

\begin{quote}
\emph{ist bei Deligne}
\end{quote}

en allemand, au milieu d'une page française.\footnote{L'allemand n'a rien
d'exceptionnel sous sa plume, mais son apparition ici, dans une note de
gestion --- « c'est chez Deligne » --- plutôt que dans une phrase mathématique,
est notable. La page ne l'explique pas.}

\pagerange{2}{2}
\section*{La seconde table : les exposés X à XXII}

Le second feuillet est entièrement de sa main et continue la numérotation. Sa
présentation change : chaque ligne s'ouvre désormais sur une lettre, D ou K, qui
tient la place que les noms occupaient au feuillet précédent.\footnote{Les deux
lettres ne sont pas développées sur la page et ne le sont pas ici. Le
rapprochement avec les deux mathématiciens dont les noms figurent ailleurs dans
le dossier --- Deligne au feuillet 1, Katz nulle part --- serait une conjecture,
et l'on s'en abstient.}

\begin{itemize}
\item \textbf{X à XII} --- intersections, intersections complètes, quadriques.
\item \textbf{XIII et XIV} --- points quadratiques ordinaires et théorème de
Lefschetz faible ; puis la formule de Milnor.\footnote{L'adjectif est peu
lisible et la transcription le marque incertain. Un point double ordinaire
--- singularité quadratique non dégénérée --- est exactement le cas où la
formule de Milnor qui suit prend sa forme la plus simple, ce qui rend la
lecture plausible sans la garantir.}
\item \textbf{XV et XVI} --- Picard--Lefschetz, deux exposés portant le même
titre et réunis par une accolade. C'est le c\oe ur quantitatif du séminaire :
la formule qui décrit exactement l'effet d'un tour de monodromie sur un cycle
évanescent.
\item \textbf{XVII} --- les pinceaux de Lefschetz, en deux moitiés soulignées
séparément : leur \emph{existence}, et leur \emph{étude cohomologique}.
\item \textbf{XIX} --- un théorème de nombres, avec Hodge--Deligne d'un côté et
Katz de l'autre.\footnote{Le numéro est corrigé : XVIII est barré et XIX écrit
à la place. Deux mots de la ligne sont illisibles.}
\item \textbf{XX} --- le théorème de Griffiths, dont le numéro est barré.
\item \textbf{XXI et XXII} --- le niveau de la cohomologie des intersections
complètes, et un calcul modulo $p$.
\end{itemize}

Sous un trait, deux lignes ferment le feuillet. La première annonce des
\emph{exposés transcendants}, avec un nom suivi d'un point d'interrogation :
Lê Dũng Tráng. La seconde n'est qu'une suite de trois symboles soulignés :
$\mathbb{F}_{q}$, $\#$, card --- un pense-bête de notation, à fixer avant que le
séminaire ne commence.

\pagerange{3}{4}
\section*{Bibliographie et rédaction}

Le troisième feuillet porte la fin d'une bibliographie : les numéros 17 à 23,
et rien d'autre. Deux grands traits de plume la barrent en diagonale.

Les sept entrées --- Griffiths, Grothendieck, Landman, Lieberman, Schmid,
Wu --- dessinent le versant transcendant du sujet : variations de structures de
Hodge, monodromie des familles de variétés complexes, thèse de Landman sur la
quasi-unipotence. Une seule correction s'y lit, et elle vaut d'être notée : au
numéro 19, en face de « S.G.A.\ on monodromy », le nom de Grothendieck est barré
et celui de Deligne écrit au-dessus.\footnote{La correction est de sa main et
c'est son propre nom qu'il raye. Elle enregistre un changement de responsabilité
sur un exposé, ce que la table du feuillet 1 laissait déjà voir ; on n'en tire
pas davantage.}

Le quatrième feuillet est un feuillet de tapuscrit, paginé (11bis) dans sa
numérotation propre et barré des mêmes traits. Il traite des \emph{bitorseurs}
et de leurs rigidifications, et c'est donc la rédaction de l'exposé VII annoncé
au premier feuillet.

L'essentiel s'y ramène à une chaîne de trois foncteurs. Étant donné deux
Groupes $P$ et $Q$ et la section unité $e_{p}$, on a

\begin{tikzcd}[column sep=large, nodes={font=\scriptsize}]
\underline{\mathrm{Ext}}(P; G) \arrow[r, "i"] & \underline{\mathrm{Bitorsrig}}(P, G) \arrow[r, "\delta"] & \underline{\mathrm{Bitorsbirig}}(P, P; G)
\end{tikzcd}

où l'on part de la catégorie des extensions de $P$ par $G$, on passe aux
bitorseurs rigidifiés, puis aux bitorseurs birigidifiés sur $P \times
P$.\footnote{Le premier foncteur est décrit sur le feuillet : à une extension
$E$ on associe le bitorseur correspondant, rigidifié par la section unité de
$E$. Le second est annoncé --- « on définit le foncteur $\delta$ par la
formule » --- et la formule manque : le feuillet s'arrête là. On ne la
restitue pas.}

Ce qui rend cette chaîne intéressante est ce qu'elle fait de la
\emph{rigidification}. Une extension porte gratuitement une section : son
élément neutre. Un bitorseur n'en porte aucune, et il faut la lui donner. Les
deux foncteurs mesurent donc exactement ce qu'on perd en oubliant la structure
de groupe, et ce qu'il faut rajouter pour la retrouver.

\pagerange{5}{5}
\section*{Les deux graphes}

Le dernier feuillet ne porte aucun mot mathématique : seulement des numéros
romains reliés par des flèches. Il y a trois graphes, et le premier est barré de
quatre grands traits croisés.\footnote{On ne le redessine pas. Ses flèches
passent sous les ratures et leurs pointes ne sont plus toutes assurées ; or un
graphe de dépendances dont une flèche manque affirme un ordre que personne n'a
écrit. Ce qu'on en voit --- une vingtaine de numéros, deux blocs, une
séparation médiane --- suffit à dire qu'il tentait déjà la même chose que les
deux suivants.}

Ce qui subsiste sous la rature est la reprise, et elle est franche : le
matériel est coupé en deux, avec une légende sous chaque moitié.

\subsection*{Théorèmes qualitatifs}

\begin{tikzcd}[column sep=small, row sep=small, nodes={font=\scriptsize}]
& & & \mathrm{I} \arrow[dlll] \arrow[ddlll] \arrow[ddll] \arrow[dddll]\\
\mathrm{II} & & \mathrm{VI} \arrow[d] \arrow[ddl] & \\
\mathrm{III} \arrow[dd] & \mathrm{IV} & \mathrm{VII} \arrow[d] \arrow[dl, dashed] & \\
& \mathrm{V} & \mathrm{VIII} \arrow[dll] & \\
\mathrm{IX} \arrow[ur] & & &
\end{tikzcd}

La forme du graphe dit son contenu. L'exposé I est une racine : quatre flèches
en partent et aucune n'y arrive. L'exposé IX est un puits : plusieurs y
arrivent. Entre les deux, deux chaînes parallèles --- II, III, IV d'un côté ;
VI, VII, VIII de l'autre --- convergent vers lui. Autrement dit, le théorème de
monodromie fournit la base, deux développements indépendants s'en séparent, et
la réduction semi-stable est ce qui les rassemble.\footnote{Une flèche du graphe
va à contre-courant : celle qui joint IX à V porte sa pointe du côté de V, alors
que toutes les autres arrivées sur IX y portent la leur. La transcription la
donne dans ce sens, qui est celui de la page, et on la conserve ici sans la
normaliser --- une dépendance retournée dans un graphe de dépendances n'est pas
un détail de dessin. Une flèche en pointillé arrive encore sur IX par la
droite, hors de ce bloc ; elle porte au-dessous un mot barré et illisible, et
n'est pas reportée.} Sous le graphe, une accolade rappelle les trois exposés
VI, VII, VIII.

\subsection*{Théorèmes quantitatifs}

\begin{tikzcd}[column sep=small, row sep=small, nodes={font=\scriptsize}]
\mathrm{XI} \arrow[d] & & \\
\mathrm{XII} \arrow[d] & & \\
\mathrm{XIII} \arrow[d] \arrow[r] & \mathrm{XIV} \arrow[r, dashed] & \mathrm{XVII} \arrow[d]\\
\mathrm{XV} \arrow[d] & & \mathrm{XVIII} \arrow[ddll]\\
\mathrm{XVI} \arrow[d] & & \\
\mathrm{XIX} & &
\end{tikzcd}

Le second graphe a une tout autre allure : c'est une chaîne, presque linéaire,
avec une seule dérivation. On descend de XI à XVI par les intersections
complètes, les quadriques et Picard--Lefschetz ; et de XIII part une branche
vers XIV, puis vers XVII et XVIII, qui rejoint la chaîne principale plus
bas.\footnote{Le X qui devrait ouvrir la chaîne est cerclé puis raturé si
lourdement qu'il n'en reste que le contour, et il est laissé hors du diagramme.
Entre XIV et XVIII court un trait terminé par une pointe et parcouru sur toute
sa longueur d'un zigzag serré, qui est la manière dont l'auteur annule une
flèche ailleurs dans le fonds ; il n'est pas transcrit non plus. Un XII posé
plus à gauche ne porte aucune flèche.}

À part, en bas à droite, quatre numéros forment une chaîne indépendante :
XIX, XX $\leftarrow$ XXI $\leftarrow$ XXII, les deux premiers cerclés.

La différence de forme entre les deux graphes est le vrai contenu de ce
feuillet. Le qualitatif est un arbre : plusieurs voies vers un même résultat. Le
quantitatif est une chaîne : un calcul, qu'il faut faire dans l'ordre.

\end{document}
